% ======================================================================
\chapter{Fundamentação teórica}
\label{cap:fundamentacao}
% ======================================================================
\instrucao{Apresente só a teoria que será usada na análise. Numere as equações, defina cada símbolo logo depois e cite a fonte. Cinco a dez páginas costumam bastar em um TCC experimental.}

\section{Ondas estacionárias em um tubo fechado}

Em um tubo fechado em uma extremidade e aberto na outra, a extremidade fechada é um nó de deslocamento e a aberta é, aproximadamente, um ventre. As ressonâncias ocorrem quando o comprimento efetivo da coluna é um múltiplo ímpar de um quarto do comprimento de onda \cite{halliday2016}:
\begin{equation}
  L_n + e = (2n-1)\,\frac{\lambda}{4}, \qquad n = 1, 2, 3, \dots
  \label{eq:ressonancia}
\end{equation}
em que $L_n$ é a posição da $n$-ésima ressonância, medida a partir da boca do tubo, e $e$ é a correção de extremidade, que leva em conta o fato de o ventre se formar um pouco fora do tubo. Para um tubo de raio $r$, a estimativa usual é $e \approx 0{,}6\,r$.

Definindo $x_n = (2n-1)/4$, a Equação~\eqref{eq:ressonancia} vira uma reta:
\begin{equation}
  L_n = \lambda\,x_n - e,
  \label{eq:reta}
\end{equation}
cujo coeficiente angular é o comprimento de onda. Com a frequência $f$ do gerador, a velocidade do som é
\begin{equation}
  v = f\,\lambda.
  \label{eq:velocidade}
\end{equation}

\section{Dependência com a temperatura}

Para um gás ideal, a velocidade do som cresce com a raiz quadrada da temperatura absoluta \cite{nussenzveig2014}. Para o ar seco, uma aproximação útil é
\begin{equation}
  v(T) = \qty{331.3}{\metre\per\second}\,\sqrt{1 + \frac{T}{\qty{273.15}{\celsius}}},
  \label{eq:temperatura}
\end{equation}
com $T$ em graus Celsius. A \qty{20}{\celsius}, ela dá cerca de \qty{343}{\metre\per\second}.

\section{Incerteza de medição}

A incerteza da velocidade vem da incerteza do comprimento de onda, já que a frequência do gerador é conhecida com incerteza relativa muito menor. Pela propagação de incertezas \cite{vuolo1996,taylor1997},
\begin{equation}
  \sigma_v = v\,\sqrt{\left(\frac{\sigma_\lambda}{\lambda}\right)^{2} + \left(\frac{\sigma_f}{f}\right)^{2}} \approx f\,\sigma_\lambda .
  \label{eq:incerteza}
\end{equation}
